\documentclass[../../main.tex]{subfiles}
\graphicspath{{\subfix{../../image/}}}

\begin{document}

\section{逼近}\label{section:ch05-逼近}

\subsection{光滑函数的局部逼近}\label{subsection:ch05-内部逼近}

\begin{theorem}[光滑函数的局部逼近定理]\label{theorem:ch05-光滑函数的局部逼近定理}
设$U\subseteq \mathbb{R}^n$是开集,固定正整数 $k$ 与 $1\leqslant p<\infty$，对每个 $\varepsilon>0$，记$U_{\varepsilon}=\{x\in U\mid \operatorname{dist}(x,\partial U)>\varepsilon\}$.再设 $u\in W^{k,p}(U)$，其中 $1\leqslant p<\infty$，令
\begin{align*}
u_{\varepsilon}\triangleq \begin{cases}
\eta _{\varepsilon}*u,&x\in U_{\varepsilon},\\
0,&x\in U\backslash U_{\varepsilon}.\\
\end{cases}
\end{align*}
则
\begin{enumerate}[(1)]
\item 对每个 $\varepsilon>0,|\alpha|\leqslant k$，有
\begin{align*}
D^{\alpha}u_{\varepsilon}=\eta_{\varepsilon}*D^{\alpha}u\,\,\text{在}\ U_{\varepsilon}\text{中},\qquad u_{\varepsilon}\in C^{\infty}(U_{\varepsilon}).
\end{align*}

\item 当 $\varepsilon\to 0$ 时，有
\begin{align*}
u_{\varepsilon}\to u \,\,\text{在}W^{k,p}_{\mathrm{loc}}(U)\text{中}.
\end{align*}
\end{enumerate}
特别地, 当 $U=\mathbb{R}^n$ 时, 由于 $\partial\mathbb{R}^n=\varnothing$, 并约定 $\operatorname{dist}(x,\varnothing)=+\infty$,则对每个 $\varepsilon>0$, 都有
\begin{align*}
U_\varepsilon=\mathbb{R}^n.
\end{align*}
此时
\begin{align*}
u_\varepsilon=\eta_\varepsilon*u
\in C^\infty(\mathbb{R}^n)\cap W^{k,p}(\mathbb{R}^n),
\end{align*}
且对每个满足 $|\alpha|\leqslant k$ 的多重指标 $\alpha$, 有
\begin{align*}
D^\alpha u_\varepsilon
=\eta_\varepsilon*D^\alpha u
\quad\text{在 }\mathbb{R}^n\text{ 中}.
\end{align*}
此外,还有
\begin{align*}
\lim\limits_{\varepsilon\to0^+}
\|u_\varepsilon-u\|_{W^{k,p}(\mathbb{R}^n)}
=0.
\end{align*}
\end{theorem}
\begin{proof}

由\hyperref[theorem:appC-光滑算子的性质]{光滑算子的性质\ref{theorem:appC-光滑算子的性质-1}}知$u_{\varepsilon}\in C^{\infty}(U_{\varepsilon})$.下面证明 :若 $|\alpha|\leqslant k$，则
\begin{align}
D^{\alpha}u_{\varepsilon}=\eta_{\varepsilon}*D^{\alpha}u,\quad\text{在}\ U_{\varepsilon}\text{中}.
\label{align:ch05-磨光导数}
\end{align}
即光滑函数 $u_{\varepsilon}$ 的通常的$\alpha$ 阶偏导数是 $u$ 的 $\alpha$ 阶弱偏导数的 $\varepsilon$-光滑化。
为证此，对固定的$x\in U_{\varepsilon}$,$\varepsilon>0$与多重指标 \(|\alpha|\leqslant k\),利用\hyperref[theorem:appC-光滑算子的性质-1]{光滑算子的性质\ref{theorem:appC-光滑算子的性质-1}}
可得
\begin{align}\label{eq::jifjiojiwjiIOJ9200--392022-1}
D^{\alpha}u_{\varepsilon}\left( x \right) = D^{\alpha}\int_U{\eta _{\varepsilon}\left( x-y \right) u\left( y \right) \,\mathrm{d}y}=\int_U{u\left( y \right) D_{x}^{\alpha}\eta _{\varepsilon}\left( x-y \right) \,\mathrm{d}y}.
\end{align}
令$F(x,y)\triangleq\eta_\varepsilon(x-y)$. 由于$\eta_\varepsilon\in C^\infty(\mathbb{R}^n)$, 故$F\in C^\infty(\mathbb{R}^{2n})$.
令$z=x-y$. 因为
\begin{align*}
\partial_{x_i}z_j=\delta_{ij},\quad\partial_{y_i}z_j=-\delta_{ij},
\end{align*}
其中$\delta_{ij}$为 Kronecker 符号, 所以对每个$1\leqslant i\leqslant n$, 多元链式法则给出
\begin{align*}
\partial_{x_i}F(x,y)
&=\sum\limits_{j=1}^{n}(\partial_j\eta_\varepsilon)(x-y)\delta_{ij}
=(\partial_i\eta_\varepsilon)(x-y),\\
\partial_{y_i}F(x,y)
&=-\sum\limits_{j=1}^{n}(\partial_j\eta_\varepsilon)(x-y)\delta_{ij}
=-(\partial_i\eta_\varepsilon)(x-y).
\end{align*}
这里$\partial_i\eta_\varepsilon$表示磨光子对第$i$个自变量的偏导数. 因而
\begin{align*}
\partial_{x_i}F=-\partial_{y_i}F.
\end{align*}
下面对$|\alpha|$归纳, 得到
\begin{align*}
D_x^\alpha F=(-1)^{|\alpha|}D_y^\alpha F.
\end{align*}
当$|\alpha|=0$时, 等式显然成立. 假设对所有$|\beta|=m$等式均成立. 对任意$|\alpha|=m+1$, 选取满足$\alpha_i\geqslant1$的$i$, 并令$\beta\triangleq\alpha-e_i$, 其中$e_i$为第$i$个坐标单位向量. 则$|\beta|=m$. 由于$F$光滑, 关于$x$和$y$的混合偏导数可以交换次序, 故由归纳假设可知
\begin{align*}
D_x^\alpha F
&=\partial_{x_i}(D_x^\beta F)=(-1)^m\partial_{x_i}(D_y^\beta F)=(-1)^mD_y^\beta(\partial_{x_i}F)\\
&=(-1)^{m+1}D_y^\beta(\partial_{y_i}F)=(-1)^{m+1}D_y^\alpha F.
\end{align*}
归纳可知, 对任意多指标$\alpha$, 都有
\begin{align}\label{eq::jifjiojiwjiIOJ9200--392022-2}
D_x^\alpha[\eta_\varepsilon(x-y)]
=(-1)^{|\alpha|}D_y^\alpha[\eta_\varepsilon(x-y)].
\end{align}
从而结合\eqref{eq::jifjiojiwjiIOJ9200--392022-1}\eqref{eq::jifjiojiwjiIOJ9200--392022-2}式可得
\begin{align*}
D^{\alpha}u_{\varepsilon}\left( x \right) = \int_Uu(y)D_x^\alpha[\eta_\varepsilon(x-y)]\,\mathrm{d}y
=(-1)^{|\alpha|}
\int_Uu(y)D_y^\alpha[\eta_\varepsilon(x-y)]\,\mathrm{d}y.
\end{align*}
现对固定的 $x\in U_{\varepsilon}$，函数 $\phi(y)\triangleq\eta_{\varepsilon}(x-y)$ 属于 $C^{\infty}_{c}(U)$。又由$\alpha$ 阶弱偏导数的定义给出
\begin{align*}
\int_{U}u(y)D_{y}^{\alpha}\eta_{\varepsilon}(x-y)\,\mathrm{d}y
=(-1)^{|\alpha|}\int_{U}\eta_{\varepsilon}(x-y)D^{\alpha}u(y)\,\mathrm{d}y.
\end{align*}
于是
\begin{align*}
D^{\alpha}u_{\varepsilon}(x)
=(-1)^{|\alpha|+|\alpha|}\int_{U}\eta_{\varepsilon}(x-y)D^{\alpha}u(y)\,\mathrm{d}y
=[\eta_{\varepsilon}*D^{\alpha}u](x).
\end{align*}
这就证明了\eqref{align:ch05-磨光导数}式。
现取开集 $V\subset \subset U$。由\eqref{align:ch05-磨光导数}式及\hyperref[theorem:appC-光滑算子的性质]{光滑算子的性质\ref{theorem:appC-光滑算子的性质-4}}，当 $\varepsilon\to 0$ 时，对每个 $|\alpha|\leqslant k$ 都有 $D^{\alpha}u_{\varepsilon}\to D^{\alpha}u$ 于 $L^{p}(V)$。从而
\begin{align*}
\|u_{\varepsilon}-u\|_{W^{k,p}(V)}^{p}
=\sum_{|\alpha|\leqslant k}\|D^{\alpha}u_{\varepsilon}-D^{\alpha}u\|_{L^{p}(V)}^{p}\to 0,\quad \varepsilon\to 0.
\end{align*}
这就证明了(2)。
\end{proof}
\begin{remark}
\hyperref[theorem:ch05-光滑函数的局部逼近定理]{光滑函数的局部逼近定理}中的局部收敛一般不能加强为
$u_\varepsilon\to u$ 于 $W^{k,p}(U)$.
虽然由\hyperref[theorem:appC-光滑算子的性质]{光滑算子的性质\ref{theorem:appC-光滑算子的性质-5}}可知
\[
\|u_\varepsilon-u\|_{L^p(U)}\to0,
\]
但整体 $L^p$ 收敛只控制函数本身的误差,
并不保证弱导数的存在及其收敛.
特别地, 在 $U\setminus U_\varepsilon$ 上将函数置零,
可能在 $\partial U_\varepsilon$ 处造成跳跃,
使得 $u_\varepsilon$ 不属于 $W^{1,p}(U)$.

例如, 取
\[
U=(0,1),\qquad u(x)\equiv1.
\]
则对任意正整数 $k$ 与 $1\leqslant p<\infty$,
均有 $u\in W^{k,p}(U)$.
当 $0<\varepsilon<\frac12$ 时,
对于 $x\in U_\varepsilon=(\varepsilon,1-\varepsilon)$,
由\hyperref[theorem:appC-光滑算子的性质]{光滑算子的性质}可得
\[
(\eta_\varepsilon*u)(x)
=\int_0^1\eta_\varepsilon(x-y)\,\mathrm{d}y=1.
\]
因此, 按照\hyperref[theorem:ch05-光滑函数的局部逼近定理]{光滑函数的局部逼近定理}的定义, 有
\[
u_\varepsilon(x)=
\begin{cases}
1,&\varepsilon<x<1-\varepsilon,\\
0,&x\in(0,\varepsilon]\cup[1-\varepsilon,1).
\end{cases}
\]
于是
\[
\|u_\varepsilon-u\|_{L^p(0,1)}^p
=\int_0^\varepsilon1\,\mathrm{d}x
+\int_{1-\varepsilon}^11\,\mathrm{d}x
=2\varepsilon\to0.
\]

下面直接利用弱导数的定义证明
$u_\varepsilon\notin W^{1,p}(0,1)$.
固定 $0<\varepsilon<\frac12$,
反设存在 $g\in L^p(0,1)$ 为 $u_\varepsilon$ 的弱导数.
则对任意 $\varphi\in C_c^\infty(0,1)$, 有
\begin{align}
\varphi(1-\varepsilon)-\varphi(\varepsilon)=\int_\varepsilon^{1-\varepsilon}\varphi'(x)\,\mathrm{d}x=\int_0^1u_\varepsilon(x)\varphi'(x)\,\mathrm{d}x
=-\int_0^1g(x)\varphi(x)\,\mathrm{d}x.\label{eq::JIOWJI3234823749235290850}
\end{align}

取\cref{definition:appC-磨光子}中 $n=1$ 时的标准磨光子 $\eta$.
由其定义可知
\[
0\leqslant\eta(t)\leqslant\eta(0)=Ce^{-1}>0,
\qquad
\eta(t)=0\quad\text{当 }|t|\geqslant1.
\]
固定 $0<\varepsilon<\frac12$.
对于
\[
0<r<\min\left\{\varepsilon,\frac{1-2\varepsilon}{2}\right\},
\]
定义
\[
\varphi _r(x)\triangleq \frac{\eta \left( \frac{x-\varepsilon}{r} \right)}{\eta (0)}.
\]
显然 $\varphi_r\in C^\infty(\mathbb R)$,
且
\[
\operatorname{spt}\varphi_r
=[\varepsilon-r,\varepsilon+r].
\]
由 $r<\varepsilon$ 以及
\[
\varepsilon+r
<\varepsilon+\frac{1-2\varepsilon}{2}
=\frac12<1,
\]
可知
\[
[\varepsilon-r,\varepsilon+r]\subset\subset(0,1).
\]
因此
$\varphi_r\in C_c^\infty(0,1)$.
此外,由$\frac{1-2\varepsilon}{r}>2>1$,且$\eta$ 在 $[-1,1]$ 外为零可得
\[
0\leqslant\varphi_r\leqslant1,\qquad
\varphi_r(\varepsilon)=1,\qquad \varphi _r(1-\varepsilon )=\frac{\eta \left( \frac{1-2\varepsilon}{r} \right)}{\eta (0)}
=0.
\]
将 $\varphi=\varphi_r$ 代入\eqref{eq::JIOWJI3234823749235290850}式, 得
\[
\int_0^1g(x)\varphi_r(x)\,\mathrm{d}x=1.
\]
由于 $\varphi_r$ 在
$[\varepsilon-r,\varepsilon+r]$ 外为零, 所以
\begin{align}
1
=\left|\int_0^1g(x)\varphi_r(x)\,\mathrm{d}x\right|
\leqslant
\int_{\varepsilon-r}^{\varepsilon+r}|g(x)|\,\mathrm{d}x.\label{eqJIOJFIOIWIOURIOWORWPPPPP}
\end{align}
然而, 因为 $g\in L^p(0,1)$ 且 $p\geqslant1$,
由\subcref{Real Analysis-proposition:Lp空间包含关系}{Real Analysis-proposition:Lp空间包含关系-1}可知 $g\in L^1(0,1)$.
再由\hyperref[Real Analysis-theorem:积分的绝对连续性]{积分的绝对连续性}知
\[
\int_{\varepsilon-r}^{\varepsilon+r}|g(x)|\,\mathrm{d}x
\to0,\qquad r\to0^+,
\]
这与\eqref{eqJIOJFIOIWIOURIOWORWPPPPP}式矛盾！
因此 $u_\varepsilon$ 不存在属于 $L^p(0,1)$ 的弱导数,
即
\[
u_\varepsilon\notin W^{1,p}(0,1).
\]
从而 $u_\varepsilon$ 也不属于 $W^{k,p}(0,1)$,
更不能在该空间中收敛到 $u$.

注意, 虽然 $u_\varepsilon$ 在两个跳跃点之外的通常导数均为零,
但这个几乎处处存在的通常导数并不是它在 $(0,1)$ 上的弱导数.
弱导数必须满足对所有测试函数成立的积分等式,
而上述反证说明不存在满足该等式的 $L^p$ 函数.

因此, 在 $U_\varepsilon$ 中成立的导数交换公式
\[
D^\alpha u_\varepsilon
=\eta_\varepsilon*D^\alpha u
\]
不能直接作为整个 $U$ 上的弱导数等式使用.
即使对每个 $D^\alpha u$ 分别应用\hyperref[theorem:appC-光滑算子的性质]{光滑算子的性质\ref{theorem:appC-光滑算子的性质-5}}, 也不能据此推出
$u_\varepsilon$ 的整体 Sobolev 收敛.

另一方面, 对任意开集 $V\subset\subset(0,1)$,
当 $\varepsilon>0$ 充分小时,
有 $V\subset(\varepsilon,1-\varepsilon)$,
因而
\[
u_\varepsilon|_V=u|_V\equiv1,
\qquad
\|u_\varepsilon-u\|_{W^{k,p}(V)}=0.
\]
所以定理中的局部收敛仍然成立.
\end{remark}
\begin{remark}
\cref{theorem:ch05-光滑函数的局部逼近定理}中的结论一般不能推广到$p=\infty$. 下面给出一个对任意正整数$k$都适用的反例.

取
\begin{align*}
U\triangleq(-1,1),\quad V\triangleq\left(-\frac12,\frac12\right),\quad
u(x)\triangleq\frac{(x_+)^k}{k!},\quad x_+\triangleq\max\{x,0\}.
\end{align*}
对任意$\varphi\in C_c^\infty(U)$及$0\leqslant j\leqslant k$, 在$(0,1)$上反复分部积分可得
\begin{align*}
\int_{-1}^{1}u(x)\varphi^{(j)}(x)\,\mathrm{d}x
=\frac{1}{k!}\int_0^1x^k\varphi^{(j)}(x)\,\mathrm{d}x
=(-1)^j\int_0^1\frac{x^{k-j}}{(k-j)!}\varphi(x)\,\mathrm{d}x.
\end{align*}
这里所有边界项均为零,这是因为:在$x=1$附近,$\varphi$及其各阶导数均为零;在$x=0$处,分部积分过程中出现的$x^k$的各阶导数的阶数至多为$k-1$,因而也均为零.所以
\begin{align*}
D^ju(x)=\frac{x^{k-j}}{(k-j)!}\chi_{(0,1)}(x),\quad 0\leqslant j\leqslant k,
\end{align*}
上述等式几乎处处成立. 这些函数均属于$L^\infty(U)$, 故$u\in W^{k,\infty}(U)$. 特别地,
\begin{align*}
D^ku(x)=
\begin{cases}
0,&-1<x<0,\\
1,&0<x<1.
\end{cases}
\end{align*}

按照定理中的方式定义$u_\varepsilon$. 当$0<\varepsilon<\frac12$时, $V\subseteq U_\varepsilon$, 因而由\hyperref[theorem:appC-光滑算子的性质]{光滑算子的性质\ref{theorem:appC-光滑算子的性质-1}}知$D^ku_\varepsilon$
在$V$上光滑, 特别是在$x=0$处连续. 对任意$\delta>0$, 由连续性, 存在$0<r<\frac12$, 使得
\begin{align*}
|D^ku_\varepsilon(x)-D^ku_\varepsilon(0)|<\delta,\quad |x|<r.
\end{align*}
因此, 在$(-r,0)$上有
\begin{align*}
|D^ku_\varepsilon(x)-D^ku(x)|
=|D^ku_\varepsilon(x)|
\geqslant |D^ku_\varepsilon(0)|-\delta,
\end{align*}
而在$(0,r)$上有
\begin{align*}
|D^ku_\varepsilon(x)-D^ku(x)|
=|D^ku_\varepsilon(x)-1|
\geqslant |D^ku_\varepsilon(0)-1|-\delta.
\end{align*}
根据本质上确界的定义,有
\begin{align*}
\|D^ku_\varepsilon-D^ku\|_{L^\infty(V)}
\geqslant\max\{|D^ku_\varepsilon(0)|,|D^ku_\varepsilon(0)-1|\}-\delta\geqslant\frac12-\delta.
\end{align*}
最后一个不等式来自
\begin{align*}
1=\left| D^ku_{\varepsilon}(0)+1-D^ku_{\varepsilon}(0) \right|\leqslant |D^ku_\varepsilon(0)|+|D^ku_\varepsilon(0)-1|
\leqslant 2\max\{|D^ku_\varepsilon(0)|,|D^ku_\varepsilon(0)-1|\}.
\end{align*}
令$\delta\to0$, 便得到对每个$0<\varepsilon<\frac12$都有
\begin{align*}
\|u_\varepsilon-u\|_{W^{k,\infty}(V)}
\geqslant\|D^ku_\varepsilon-D^ku\|_{L^\infty(V)}
\geqslant\frac12.
\end{align*}
因此, 当$\varepsilon\to0$时, $u_\varepsilon$不在$W^{k,\infty}(V)$中收敛到$u$, 从而也不在$W_{\mathrm{loc}}^{k,\infty}(U)$中收敛到$u$.

直观上, 这个反例的关键障碍来自最高阶弱导数的跳跃. 函数$D^ku$在原点处从$0$跳到$1$, 而$D^ku_\varepsilon$是连续函数, 必须通过一段连续过渡来连接两侧的值. 随着$\varepsilon\to0$, 过渡区域可以越来越窄, 但其中的误差始终无法一致地变小.

具体而言, 对通常的非负、积分为$1$且支集包含于$[-\varepsilon,\varepsilon]$的磨光子, 由\eqref{align:ch05-磨光导数}可知, 在$V$内有
\begin{align*}
0\leqslant D^ku_\varepsilon\leqslant1,\quad
D^ku_\varepsilon(x)=D^ku(x)\quad\text{当 }|x|>\varepsilon.
\end{align*}
于是对每个$1\leqslant p<\infty$,有
\begin{align*}
\|D^ku_\varepsilon-D^ku\|_{L^p(V)}^p
=\int_{-\varepsilon}^{\varepsilon}
|D^ku_\varepsilon(x)-D^ku(x)|^p\,\mathrm{d}x
\leqslant2\varepsilon\to0.
\end{align*}
这说明有限$p$的积分范数可以利用误差区域的测度趋于零而趋于零, 但$L^\infty$范数只取误差的本质上确界, 无法利用区域变窄来消除误差. 因而, \cref{theorem:ch05-光滑函数的局部逼近定理}中的光滑性及导数交换公式在$p=\infty$时仍然成立, 不成立的是最后使用的磨光逼近在$L^p$范数下的收敛结论.
\end{remark}

\subsection{光滑函数的整体逼近}\label{subsection:ch05-光滑函数逼近}

\begin{lemma}[光滑单位分解]\label{lemma:ch05-光滑单位分解-p7R3}
设 $U\subseteq\mathbb{R}^{n}$ 为开集, $\{V_i\}_{i=0}^{\infty}$ 是 $U$ 的可数开覆盖, 且 $V_i\subset\subset U$. 假设该覆盖在 $U$ 中局部有限, 即每个 $x\in U$ 都有一个邻域, 仅与有限多个 $V_i$ 相交. 则存在 $\zeta_i\in C_c^\infty(V_i)$, 使得
\begin{align*}
0\leqslant\zeta_i\leqslant1,\quad
\sum\limits_{i=0}^{\infty}\zeta_i(x)=1,\quad x\in U.
\end{align*}
这里将每个 $\zeta_i$ 在 $V_i$ 外延拓为零. 此外, 每个紧集 $K\subseteq U$ 只与有限多个 $\operatorname{spt}\zeta_i$ 相交, 即上述和在 $U$ 中为局部有限和.
\end{lemma}
\begin{proof}

若 $U=\varnothing$, 取所有 $\zeta_i=0$ 即可. 以下设 $U\neq\varnothing$.

\noindent$\mathbf{Step}\mathbf{1}.${\heiti 构造从属于开覆盖的紧集覆盖.}

对每个 $i\geqslant0$, 由\subcref{proposition:点集间的距离的性质}{proposition:点集间的距离的性质-3},定义
\begin{align*}
a_i(x)\triangleq\operatorname{dist}(x,\mathbb{R}^{n}\setminus V_i),\quad x\in\mathbb{R}^{n}.
\end{align*}
由于 $V_i$ 是有界开集, $\mathbb{R}^{n}\setminus V_i$ 是非空闭集, 故 $a_i$ 是连续函数, 并且
\begin{align*}
a_i(x)>0\Longleftrightarrow x\in V_i.
\end{align*}
对 $x\in U$, 令
\begin{align}\label{eq::JIJIOWwqqqwqqqwqq}
A(x)\triangleq\max\limits_{i\geqslant0}a_i(x).
\end{align}
下面说明这个最大值确实存在, 且 $A$ 是 $U$ 上的正连续函数.

任取 $x\in U$. 由覆盖的局部有限性, 存在包含 $x$ 的开集 $W_x\subseteq U$, 使得指标集
\begin{align*}
I_x\triangleq\{i\geqslant0:W_x\cap V_i\neq\varnothing\}
\end{align*}
为有限集. 当 $i\notin I_x$ 时, $a_i$ 在 $W_x$ 上恒为零. 又因为 $\{V_i\}$ 覆盖 $U$, 所以每个 $y\in W_x$ 至少属于一个 $V_i$, 从而
\begin{align*}
A(y)=\max\limits_{i\in I_x}a_i(y)>0,\quad y\in W_x.
\end{align*}
因此, $A$ 在每一点附近(即$W_x$)上是有限个连续函数的最大值, 故连续. 由 $x$ 的任意性, 可知 $A$ 在 $U$ 上连续且处处为正.

对每个 $i\geqslant0$, 定义
\begin{align*}
K_i\triangleq\left\{x\in U:a_i(x)\geqslant\frac{A(x)}{2}\right\}.
\end{align*}
若 $x\in K_i$, 则 $a_i(x)\geqslant\frac{A(x)}{2}>0$, 因而 $x\in V_i$. 所以 $K_i\subseteq V_i$.

由 $a_i$ 和 $A$ 的连续性, $K_i$ 在 $U$ 中相对闭. 又因为 $\overline{V_i}$ 是包含于 $U$ 的紧集, 且 $K_i\subseteq\overline{V_i}$, 所以 $K_i$ 是 $\overline{V_i}$ 的闭子集, 从而 $K_i$ 也是紧集.

最后, 对任意 $x\in U$,由\eqref{eq::JIJIOWwqqqwqqqwqq}式中$A(x)$的定义知,可选取一个满足 $a_i(x)=A(x)$ 的指标 $i$, 即有 $x\in K_i$. 因此
\begin{align}\label{eq::JIOJIOJEWWQWERTAAPPP}
U=\bigcup\limits_{i=0}^{\infty}K_i.
\end{align}

\noindent$\mathbf{Step}\mathbf{2}.${\heiti 构造在各个紧集$K_i$上等于$1$的光滑截断函数.}

取非负标准光滑化子 $\eta\in C_c^\infty(B(0,1))$, 满足
\begin{align*}
\int_{\mathbb{R}^{n}}\eta(x)\,\mathrm{d}x=1,
\end{align*}
并令
\begin{align*}
\eta_r(x)\triangleq\frac{1}{r^{n}}\eta\left(\frac{x}{r}\right),\quad r>0.
\end{align*}
若 $K_i=\varnothing$, 取 $\varphi_i=0$. 若 $K_i\neq\varnothing$, 由 $K_i$ 紧且 $K_i\subseteq V_i$, 结合\subcref{proposition:点集间的距离的性质}{proposition:点集间的距离的性质-3}可知
\begin{align*}
b_i\triangleq\operatorname{dist}(K_i,\mathbb{R}^{n}\setminus V_i)>0.
\end{align*}
令
\begin{align*}
r_i\triangleq\frac{b_i}{4},\quad
E_i\triangleq\{y\in\mathbb{R}^{n}:\operatorname{dist}(y,K_i)<2r_i\},
\end{align*}
并定义
\begin{align*}
\varphi_i(x)\triangleq(\eta_{r_i}*\chi_{E_i})(x)
=\int_{\mathbb{R}^{n}}\eta_{r_i}(z)\chi_{E_i}(x-z)\,\mathrm{d}z.
\end{align*}
由于 $E_i$ 是有界开集, $\chi_{E_i}\in L^1(\mathbb{R}^{n})$. 因而由\hyperref[theorem:appC-光滑算子的性质]{光滑算子的性质\ref{theorem:appC-光滑算子的性质-1}}, 有 $\varphi_i\in C^\infty(\mathbb{R}^{n})$. 又因为 $\eta_{r_i}\geqslant0$, 且其积分等于 $1$, 所以
\begin{align*}
0\leqslant\varphi_i(x)\leqslant1,\quad x\in\mathbb{R}^{n}.
\end{align*}
若 $\operatorname{dist}(x,K_i)<r_i$, 则对每个 $z\in\operatorname{spt}\eta_{r_i}$, 由 $|z|\leqslant r_i$ 可得
\begin{align*}
\operatorname{dist}(x-z,K_i)
\leqslant\operatorname{dist}(x,K_i)+|z|
<2r_i.
\end{align*}
因此 $x-z\in E_i$, 从而
\begin{align*}
\varphi_i(x)
=\int_{\mathbb{R}^{n}}\eta_{r_i}(z)\,\mathrm{d}z
=1.
\end{align*}
这说明 $\varphi_i$ 在开集 $\{x:\operatorname{dist}(x,K_i)<r_i\}$ 上恒等于 $1$, 特别地, 在 $K_i$ 上恒等于 $1$.

另一方面, 由\hyperref[theorem:appC-光滑算子的性质-6]{光滑算子的性质\ref{theorem:appC-光滑算子的性质-6}}, 有
\begin{align*}
\operatorname{spt}\varphi_i
\subseteq\overline{E_i}+\overline{B(0,r_i)}
\subseteq\{x\in\mathbb{R}^{n}:\operatorname{dist}(x,K_i)\leqslant3r_i\}.
\end{align*}
右端是有界闭集, 且由 $3r_i<b_i$ 可知它包含于 $V_i$. 所以 $\operatorname{spt}\varphi_i$ 是包含于 $V_i$ 的紧集, 即 $\varphi_i\in C_c^\infty(V_i)$.

\noindent$\mathbf{Step}\mathbf{3}.${\heiti 证明归一化分母$\Phi(x)$是正的光滑函数.}

对 $x\in U$, 定义
\begin{align*}
\Phi(x)\triangleq\sum\limits_{i=0}^{\infty}\varphi_i(x).
\end{align*}
任取 $x\in U$, 仍取$\mathbf{Step}\mathbf{1}$中的开邻域 $W_x$ 及有限指标集 $I_x$. 因为 $\operatorname{spt}\varphi_i\subseteq V_i$, 当 $i\notin I_x$ 时, $\varphi_i$ 在 $W_x$ 上恒为零. 因此
\begin{align}\label{eq::JIJIWIOIOWIWIWIWWIOWIWUIQIWIQUI21213}
\Phi(y)=\sum\limits_{i\in I_x}\varphi_i(y),\quad y\in W_x.
\end{align}
所以 $\Phi$ 在每一点附近(即$W_x$)都是有限个光滑函数之和, 即$\Phi$在$U$上是局部有限的.从而 $\Phi\in C^\infty(U)$.

由\eqref{eq::JIOJIOJEWWQWERTAAPPP}式, 每个 $x\in U$ 至少属于一个 $K_i$. 对这个指标 $i$,由$\mathbf{Step}\mathbf{2}$有 $\varphi_i(x)=1$. 再由所有 $\varphi_i$ 均非负, 得
\begin{align*}
\Phi(x)\geqslant1,\quad x\in U.
\end{align*}

\noindent$\mathbf{Step}\mathbf{4}.${\heiti 归一化并验证全部结论.}

对每个 $i\geqslant0$, 定义
\begin{align*}
\zeta_i(x)\triangleq\frac{\varphi_i(x)}{\Phi(x)},\quad x\in U.
\end{align*}
由于 $\Phi$ 光滑且处处为正, $\zeta_i\in C^\infty(U)$. 又因为$\operatorname{spt}\zeta_i
\subseteq\operatorname{spt}\varphi_i
\subset\subset V_i,$
因此 $\zeta_i\in C_c^\infty(V_i)$.

由$0\leqslant\varphi_i\leqslant\Phi$, 有 $0\leqslant\zeta_i\leqslant1$. 同时, 利用$\mathbf{Step}\mathbf{3}$已经证明的$\Phi$在$U$上是局部有限的,即由\eqref{eq::JIJIWIOIOWIWIWIWWIOWIWUIQIWIQUI21213}式可得
\begin{align*}
\sum_{i=0}^{\infty}{\zeta _i\left( x \right)}=\frac{1}{\Phi \left( x \right)}\sum_{i=0}^{\infty}{\varphi _i\left( x \right)}=\frac{\sum\limits_{x\in I_x}{\varphi _i\left( x \right)}}{\sum\limits_{x\in I_x}{\varphi _i\left( x \right)}}=1,\quad x\in U.
\end{align*}

最后, 任取紧集 $K\subseteq U$. 对每个 $x\in K$, 仍取$\mathbf{Step}\mathbf{1}$中的一个只与有限多个 $V_i$ 相交的开邻域 $W_x\subseteq U$. 由$K$的紧性, 存在有限多个这样的邻域, 使得
\begin{align*}
K\subseteq\bigcup\limits_{\ell=1}^{m}W_{x_\ell}.
\end{align*}
每个 $W_{x_\ell}$ 只与有限多个 $V_i$ 相交, 而 $\operatorname{spt}\zeta_i\subseteq V_i$, 因此 $K$ 也只与有限多个 $\operatorname{spt}\zeta_i$ 相交.
\end{proof}

\begin{corollary}
设 $K\subseteq\mathbb{R}^{n}$ 为紧集, $V_0,\cdots,V_N$ 为 $\mathbb{R}^{n}$ 中的开集, 且
\begin{align*}
K\subseteq\bigcup\limits_{i=0}^{N}V_i.
\end{align*}
则存在 $\zeta_i\in C_c^\infty(V_i)$, 使得
\begin{align*}
0\leqslant\zeta_i\leqslant1,\quad
\sum\limits_{i=0}^{N}\zeta_i(x)=1,\quad x\in K.
\end{align*}
\end{corollary}
\begin{remark}
事实上, 可以使这些函数之和在 $K$ 的某个开邻域上恒等于 $1$.
\end{remark}
\begin{proof}

若 $K=\varnothing$, 取所有 $\zeta_i=0$ 即可. 以下设 $K\neq\varnothing$.

\noindent$\mathbf{Step}\mathbf{1}.${\heiti 将有限开覆盖扩充为整个空间的局部有限开覆盖.}

由 $K$ 的有界性, 可取 $R>0$, 使得 $K\subseteq B(0,R)$. 定义
\begin{gather*}
W_i\triangleq V_i\cap B(0,2R),\,\,i=0,\cdots,N,\quad W_{N+1}\triangleq B(0,2R)\setminus K,
\\
W_{N+1+m}\triangleq
\{x\in\mathbb{R}^{n}:mR<|x|<(m+2)R\},\,\, m=1,2,\cdots.
\end{gather*}
由于 $K$ 是闭集, 上述每个 $W_j$ 都是有界开集, 因而 $W_j\subset\subset\mathbb{R}^{n}$.

由 $V_0,\cdots,V_N$ 覆盖 $K$, 有
\begin{align*}
\bigcup\limits_{i=0}^{N+1}W_i=B(0,2R).
\end{align*}
另一方面,我们有
\begin{align*}
\bigcup\limits_{m=1}^{\infty}W_{N+1+m}
=\{x\in\mathbb{R}^{n}:|x|>R\}.
\end{align*}
因此 $\{W_j\}_{j=0}^{\infty}$ 覆盖 $\mathbb{R}^{n}$.

注意到任意有界球只与有限多个$\{W_{N+1+m}\}_{m=1}^{\infty}$中集合相交, 再加上 $W_0,\cdots,W_{N+1}$ 只有有限个, 可知这个覆盖在 $\mathbb{R}^{n}$ 中局部有限. 此外, 还有
\begin{align*}
W_i\subseteq V_i,\quad i=0,\cdots,N,\qquad
W_j\cap K=\varnothing,\quad j\geqslant N+1.
\end{align*}

\noindent$\mathbf{Step}\mathbf{2}.${\heiti 应用\cref{lemma:ch05-光滑单位分解-p7R3}并取出所需的有限个函数.}

将\cref{lemma:ch05-光滑单位分解-p7R3}应用于开集 $\mathbb{R}^{n}$ 及其覆盖 $\{W_j\}_{j=0}^{\infty}$, 得到 $\psi_j\in C_c^\infty(W_j)$, 满足
\begin{align*}
0\leqslant\psi_j\leqslant1,\quad
\sum\limits_{j=0}^{\infty}\psi_j(x)=1,\quad x\in\mathbb{R}^{n}.
\end{align*}
令
\begin{align*}
\zeta_i\triangleq\psi_i,\quad i=0,\cdots,N.
\end{align*}
因为 $W_i\subseteq V_i$, 所以 $\zeta_i\in C_c^\infty(V_i)$, 且 $0\leqslant\zeta_i\leqslant1$.

\noindent$\mathbf{Step}\mathbf{3}.${\heiti 证明这些函数之和在 $K$ 上恒等于$1$.}

由\noindent$\mathbf{Step}\mathbf{2}$知$\psi_j\in C_c^\infty(W_j)\,(j=0,1,\cdots)$,进而
\begin{gather*}
\operatorname{spt}\psi_{N+1}
\subseteq W_{N+1}
=B(0,2R)\setminus K,\\
\operatorname{spt}\psi_{N+1+m}
\subseteq W_{N+1+m}
\subseteq\{y\in\mathbb{R}^{n}:|y|>R\}.
\end{gather*}
任取$x\in K\subseteq B(0,R)$,于是对每个 $m\geqslant1$,有$\psi_{N+1+m}(x)=0$.因此
\begin{align*}
\sum\limits_{i=0}^{N}\zeta_i(x)
=\sum\limits_{j=0}^{\infty}\psi_j(x)
=1,\quad \forall x\in K.
\end{align*}
\end{proof}

\begin{theorem}[光滑函数的整体逼近定理]\label{theorem:ch05-光滑函数的整体逼近定理}
设 $U\subset \mathbb{R}^n$是有界开集，并设 $u\in W^{k,p}(U)$，其中 $k\in \mathbb{N},1\leqslant p<\infty$。则存在函数 $u_{m}\in C^{\infty}(U)\cap W^{k,p}(U)$ 使
\begin{align*}
u_{m}\to u,\quad\text{在}\ W^{k,p}(U)\text{中}.
\end{align*}
\end{theorem}
\begin{proof}
若 $U=\varnothing$, 结论显然成立. 以下设 $U\neq\varnothing$.

\noindent$\mathbf{Step}\mathbf{1}.${\heiti 构造局部有限的开覆盖及从属于它的光滑单位分解.}

记 $d(x)\triangleq\operatorname{dist}(x,\partial U)$, 并定义
\begin{align*}
U_i\triangleq\left\{x\in U:d(x)>\frac{1}{i}\right\},
\quad i=1,2,\cdots.
\end{align*}
由距离函数的连续性可知, 各个 $U_i$ 都是开集, 且 $U_i\subseteq U_{i+1}$. 又因为每个 $x\in U$ 都满足 $d(x)>0$, 所以
\begin{align}
U=\bigcup\limits_{i=1}^{\infty}U_i.
\label{eqLLLLLLLLPMMMMMOJOOKMQQQAAA}
\end{align}
若 $x\in\overline{U_i}$, 则 $x\in\overline{U}$ 且 $d(x)\geqslant\frac{1}{i}>0$.
由\subcref{Real Analysis-proposition:点集间的距离的性质}{Real Analysis-proposition:点集间的距离的性质-2}, 有 $x\notin\partial U$, 从而 $x\in U$. 因此
\begin{align*}
\overline{U_i}
\subseteq\left\{x\in U:d(x)\geqslant\frac{1}{i}\right\}
\subseteq U_{i+1}.
\end{align*}
由于 $U$ 有界, $\overline{U_i}$ 是紧集, 故 $U_i\subset \subset U$.

定义
\begin{align*}
V_0\triangleq U_3,\quad W_0\triangleq U_4,
\end{align*}
并对 $i\geqslant1$ 定义
\begin{align*}
V_i\triangleq U_{i+3}\setminus\overline{U_{i+1}},
\quad
W_i\triangleq U_{i+4}\setminus\overline{U_i}.
\end{align*}
这些集合都是开集, 且 $\overline{V_0}\subseteq W_0$. 对每个 $i\geqslant1$, 有
\begin{align*}
\overline{V_i}
\subseteq
\left\{x\in U:
\frac{1}{i+3}\leqslant d(x)\leqslant\frac{1}{i+1}\right\}
\subseteq U_{i+4}\setminus\overline{U_i}
=W_i.
\end{align*}
再由 $\overline{W_i}\subseteq\overline{U_{i+4}}$, 可得
\begin{align*}
V_i\subset \subset W_i\subset \subset U,\quad i=0,1,\cdots.
\end{align*}

下面验证 $\{V_i\}_{i=0}^{\infty}$ 覆盖 $U$.
任取 $x\in U$. 若 $x\in U_3$, 则 $x\in V_0$.
若 $x\notin U_3$, 则由\eqref{eqLLLLLLLLPMMMMMOJOOKMQQQAAA}及\hyperref[Topology(Munkres)-theorem:ch01-良序性质]{良序性质}, 可取使 $x\in U_j$ 成立的最小正整数 $j$. 此时 $j\geqslant4$, 且
\begin{align*}
\frac{1}{j}<d(x)\leqslant\frac{1}{j-1}<\frac{1}{j-2}.
\end{align*}
因此 $x\in U_j$ 且 $x\notin\overline{U_{j-2}}$, 从而
\begin{align*}
x\in U_j\setminus\overline{U_{j-2}}=V_{j-3}.
\end{align*}
故 $\{V_i\}_{i=0}^{\infty}$ 确实覆盖 $U$.

再验证局部有限性.
任取非空紧集 $K\subseteq U$.
由\subcref{Real Analysis-proposition:点集间的距离的性质}{Real Analysis-proposition:点集间的距离的性质-3}, 有
\begin{align*}
\operatorname{dist}(K,\partial U)>0.
\end{align*}
因而可选取正整数 $N$, 使得 $K\subseteq U_N$.
对每个 $i\geqslant N$, 有 $K\subseteq U_i$, 而
\begin{align*}
W_i\subseteq U\setminus\overline{U_i},
\end{align*}
所以 $K\cap W_i=\varnothing$.
这说明每个紧子集 $K\subseteq U$ 只与有限多个 $W_i$ 相交.
在 $U$ 中每一点附近取一个闭包包含于 $U$ 的开球$B$,从而$B$也只与有限多个$W_i$相交,即知 $\{W_i\}_{i=0}^{\infty}$ 在 $U$ 中局部有限.
由 $V_i\subseteq W_i$知$\{V_i\}_{i=0}^{\infty}$ 也在 $U$ 中局部有限.

由\hyperref[lemma:ch05-光滑单位分解-p7R3]{光滑单位分解引理}, 存在从属于 $\{V_i\}_{i=0}^{\infty}$ 的光滑单位分解 $\{\zeta_i\}_{i=0}^{\infty}$, 满足
\begin{align}
0\leqslant\zeta_i\leqslant1,\quad
\zeta_i\in C_c^\infty(V_i),\quad
\sum\limits_{i=0}^{\infty}\zeta_i=1
\quad\text{在 }U\text{中}.
\label{align:ch05-单位分解-a8L2}
\end{align}
将各个$\zeta_i$零延拓到$U$后仍是$U$上具有紧支集的光滑函数.
从而
\begin{align}\label{eqJIOJFIOIWIOURIOWORWPPPPPJAAFAAFAAWFASCcc2}
\operatorname{spt}\zeta_i\subseteq V_i\subseteq W_i\subset \subset U,
\end{align}
且$\{W_i\}_{i=0}^{\infty}$ 在 $U$ 中局部有限,故上述单位分解之和在 $U$ 中局部有限.

\noindent$\mathbf{Step}\mathbf{2}.${\heiti 分别磨光各个乘积, 并同时控制支集与逼近误差.}

取标准磨光核 $\eta\in C_c^\infty(B(0,1))$, 满足
\begin{align*}
\eta\geqslant0,\quad
\int_{\mathbb{R}^n}\eta(x)\,\mathrm{d}x=1,
\end{align*}
并记
\begin{align*}
\eta_\varepsilon(x)
\triangleq\frac{1}{\varepsilon^n}
\eta\left(\frac{x}{\varepsilon}\right),
\quad \varepsilon>0.
\end{align*}

固定一个 $i\geqslant0$.
由\hyperref[theorem:ch05-弱导数性质-4]{弱导数性质\ref{theorem:ch05-弱导数性质-4}}, 有
\begin{align}\label{eqJIOJFIOIWIOURIOWORWPPPPPJAAFAAFAAWFASCcc1}
\zeta_i u\in W^{k,p}(U).
\end{align}
并且, 对每个满足 $|\alpha|\leqslant k$ 的多重指标 $\alpha$,有
\begin{align}\label{eqJIOJFIOIWIOURIOWORWPPPPPJAAFAAFAAWFASCcc}
D^\alpha(\zeta_i u)
=\sum\limits_{\beta\leqslant\alpha}
\binom{\alpha}{\beta}
D^\beta\zeta_i\,D^{\alpha-\beta}u.
\end{align}
再结合\eqref{eqJIOJFIOIWIOURIOWORWPPPPPJAAFAAFAAWFASCcc2}式知
\begin{align}\label{HJIOJnvzcvzn}
\operatorname{spt}D^\alpha(\zeta_i u)\subseteq \operatorname{spt}\zeta_i\subseteq V_i\subseteq W_i\subset \subset U.
\end{align}

先考虑 $\zeta_i\not\equiv0$ 的情形, 由\eqref{align:ch05-单位分解-a8L2}式知$\operatorname{spt}\zeta_i$ 是包含于 $W_i$ 的非空紧集,再由\subcref{Real Analysis-proposition:点集间的距离的性质}{Real Analysis-proposition:点集间的距离的性质-3}, 可定义两个正数
\begin{align*}
\rho_i
\triangleq\operatorname{dist}
\bigl(\operatorname{spt}\zeta_i,\mathbb{R}^n\setminus W_i\bigr)>0,\quad
r_i
\triangleq\operatorname{dist}
\bigl(\overline{W_i},\partial U\bigr)>0.
\end{align*}
再令
\begin{align*}
a_i\triangleq\frac{1}{2}\min\{\rho_i,r_i\}>0.
\end{align*}
则对每个 $0<\varepsilon<a_i$, 都有
\begin{align}\label{eq::JIJOWNVNZNJCS}
\operatorname{spt}\zeta_i+\overline{B(0,\varepsilon)}\subseteq W_i,
\quad
\overline{W_i}\subseteq
U_\varepsilon
\triangleq\{x\in U:d(x)>\varepsilon\}.
\end{align}
事实上, 若 $y=x+z\in \operatorname{spt}\zeta_i+\overline{B(0,\varepsilon)}$, 其中 $x\in \operatorname{spt}\zeta_i$ 且 $|z|\leqslant\varepsilon$, 则对任意 $w\in\mathbb{R}^n\setminus W_i$, 由三角不等式及 $\rho_i$ 的定义, 有
\begin{align*}
|y-w|
=|(x-w)+z|
\geqslant|x-w|-|z|
\geqslant\rho_i-|z|.
\end{align*}
对 $w\in\mathbb{R}^n\setminus W_i$ 取下确界, 再利用 $|z|\leqslant\varepsilon<\rho_i$,即得
\begin{align*}
\operatorname{dist}
\bigl(y,\mathbb{R}^n\setminus W_i\bigr)
\geqslant\rho_i-|z|
\geqslant\rho_i-\varepsilon>0,
\end{align*}
由\subcref{Real Analysis-proposition:点集间的距离的性质}{Real Analysis-proposition:点集间的距离的性质-2}知$\{y\}\cap (\mathbb{R}^n\setminus W_i)=\varnothing $,故$y\in W_i$.
另一方面, 对每个 $x\in\overline{W_i}$, 有
\begin{align*}
d(x)\geqslant r_i>\varepsilon,
\end{align*}
故 $x\in U_\varepsilon$.

将 $\zeta_i u$ 零延拓到 $\mathbb{R}^n$, 记为
\begin{align*}
f_i(x)\triangleq
\begin{cases}
\zeta_i(x)u(x),&x\in U,\\
0,&x\in\mathbb{R}^n\setminus U.
\end{cases}
\end{align*}
则由\eqref{eqJIOJFIOIWIOURIOWORWPPPPPJAAFAAFAAWFASCcc2}\eqref{eqJIOJFIOIWIOURIOWORWPPPPPJAAFAAFAAWFASCcc1}\eqref{HJIOJnvzcvzn}式知
\begin{align*}
\left\| f_i \right\| _{W^{k,p}\left( \mathbb{R} ^n \right)}=\left\| \zeta _iu \right\| _{W^{k,p}\left( \mathbb{R} ^n \right)}=\left\| \zeta _iu \right\| _{W^{k,p}\left( U \right)}<\infty ,
\end{align*}
因此$f_i\in W^{k,p}(\mathbb{R}^n)$.
对 $\varepsilon>0$, 定义
\begin{align*}
u_{i,\varepsilon}(x)
\triangleq(\eta_\varepsilon*f_i)(x)
=\int_{\mathbb{R}^n}
\eta_\varepsilon(x-y)f_i(y)\,\mathrm{d}y,
\quad x\in\mathbb{R}^n.
\end{align*}
由\hyperref[theorem:appC-光滑算子的性质-1]{光滑算子的性质\ref{theorem:appC-光滑算子的性质-1}}知$u_{i,\varepsilon}\in C^\infty(\mathbb{R}^n)$.
又由\hyperref[theorem:appC-光滑算子的性质-6]{光滑算子的性质\ref{theorem:appC-光滑算子的性质-6}}和\eqref{eq::JIJOWNVNZNJCS}式, 对每个 $0<\varepsilon<a_i$, 有
\begin{align}
\operatorname{spt}u_{i,\varepsilon}
\subseteq \operatorname{spt}\zeta_i+\overline{B(0,\varepsilon)}
\subseteq W_i.
\label{eq::jiowjioWJioFJIOUWE-12132}
\end{align}
由于 $\operatorname{spt}\zeta_i+\overline{B(0,\varepsilon)}$ 是紧集, 得
\begin{align}\label{eqHuEhiojwije}
u_{i,\varepsilon}\in C_c^\infty(W_i),
\quad 0<\varepsilon<a_i.
\end{align}

由于 $W_i\subset \subset U\subset \subset \mathbb{R}^n$, 故利用\hyperref[theorem:ch05-光滑函数的局部逼近定理]{光滑函数的局部逼近定理}, 可得
\begin{align*}
\lim_{\varepsilon \rightarrow 0^+} \left\| u_{i,\varepsilon}-f_i \right\| _{W^{k,p}(W_i)}=\lim_{\varepsilon \rightarrow 0^+} \left\| u_{i,\varepsilon}-\zeta _iu \right\| _{W^{k,p}(W_i)}=0.
\end{align*}
对于 $0<\varepsilon<a_i$,由\eqref{eqJIOJFIOIWIOURIOWORWPPPPPJAAFAAFAAWFASCcc2}\eqref{eqJIOJFIOIWIOURIOWORWPPPPPJAAFAAFAAWFASCcc}\eqref{eqHuEhiojwije}式知,函数 $u_{i,\varepsilon}$ 的各阶导数以及 $\zeta_i u$ 的 $k$ 阶以内弱导数都在 $U\setminus W_i$ 上几乎处处为零, 所以
\begin{align*}
\|u_{i,\varepsilon}-\zeta_i u\|_{W^{k,p}(U)}
=
\|u_{i,\varepsilon}-\zeta_i u\|_{W^{k,p}(W_i)}.
\end{align*}
因此
\begin{align*}
\lim\limits_{\varepsilon\to0^+}
\|u_{i,\varepsilon}-\zeta_i u\|_{W^{k,p}(U)}
=0.
\end{align*}
从而对$\forall \delta>0$,存在 $b_i>0$, 使得只要$0<\varepsilon<\min\{a_i,b_i\},$就有
\begin{align}\label{eq:JIFJIOUWINOUNIWNIONIOMXIZZ}
\|u_{i,\varepsilon}-\zeta_i u\|_{W^{k,p}(U)}
<\frac{\delta}{2^{i+1}}.
\end{align}
现在取
\begin{align*}
\varepsilon_i
\triangleq\frac{1}{2}\min\{a_i,b_i\},
\quad
u_i\triangleq u_{i,\varepsilon_i}.
\end{align*}

再考虑$\zeta_i\equiv0$的情形,直接取 $\varepsilon_i=1$ 及 $u_i\equiv 0$即可.于是由\eqref{eqHuEhiojwije}\eqref{eq:JIFJIOUWINOUNIWNIONIOMXIZZ}式,对所有 $i\geqslant0$, 均有
\begin{align}
u_i\in C_c^\infty(W_i),\quad
\|u_i-\zeta_i u\|_{W^{k,p}(U)}
\leqslant\frac{\delta}{2^{i+1}}.
\label{align:ch05-逼近误差-e6Q2}
\end{align}

\noindent$\mathbf{Step}\mathbf{3}.${\heiti 利用误差级数的收敛性完成求和, 并构造逼近序列.}

令
\begin{align*}
v\triangleq\sum\limits_{i=0}^{\infty}u_i.
\end{align*}
任取 $x\in U$. 由 $\{W_i\}_{i=0}^{\infty}$ 的局部有限性, 存在包含 $x$ 的开集 $W_x\subseteq U$, 使得指标集
\begin{align*}
I_x\triangleq\{i\geqslant0:W_x\cap W_i\neq\varnothing\}
\end{align*}
为有限集.
由\eqref{align:ch05-逼近误差-e6Q2}式, 有 $\operatorname{spt}u_i\subseteq W_i$, 所以当 $i\notin I_x$ 时, $u_i$ 在 $W_x$ 上恒为零.
因此
\begin{align*}
v(y)=\sum\limits_{i\in I_x}u_i(y),\quad y\in W_x.
\end{align*}
右端是有限个光滑函数之和. 由 $x$ 的任意性, 可知 $v$ 是良定义的, 且 $v\in C^\infty(U)$.

由$\sum\limits_{i=0}^{\infty}\zeta_i=1$知$u=\sum_{i=0}^{\infty}\zeta_{i}u$,从而
\begin{align*}
\|v-u\|_{W^{k,p}(U)}
=\left\|\sum\limits_{i=0}^{\infty}(u_i-\zeta_i u)\right\|_{W^{k,p}(U)}
\leqslant\sum\limits_{i=0}^{\infty}
\|u_i-\zeta_i u\|_{W^{k,p}(U)}
\leqslant\delta\sum\limits_{i=0}^{\infty}\frac{1}{2^{i+1}}
=\delta.
\end{align*}
进而$v-u\in W^{k,p}(U)$,结合 $u\in W^{k,p}(U)$和$W^{k,p}(U)$空间的线性性, 得 $v\in C^\infty(U)\cap W^{k,p}(U)$.

最后, 对每个正整数 $m$, 取 $\delta=\frac{1}{m}$, 并将对应的 $v$ 记为 $u_m$, 则
\begin{align*}
u_m\in C^\infty(U)\cap W^{k,p}(U),\quad
\|u_m-u\|_{W^{k,p}(U)}
\leqslant\frac{1}{m}\longrightarrow0.
\end{align*}
\end{proof}



\subsection{直到边界光滑的函数整体逼近}\label{subsection:ch05-边界逼近}

\begin{lemma}[正交平移变换为等距双射]\label{lemma:正交平移变换为等距双射}
设 $Q$ 为 $n$ 阶正交矩阵，$x_0\in\mathbb{R}^n$，定义
\begin{align*}
T(y)\triangleq x_0+Q^T y,\qquad y\in\mathbb{R}^n.
\end{align*}
则 $T$ 是 $\mathbb{R}^n$ 到自身的等距双射，其逆映射为
\begin{align*}
T^{-1}(x)=Q(x-x_0),\qquad x\in\mathbb{R}^n.
\end{align*}
并且 $T$ 与 $T^{-1}$ 都是 Lipschitz 映射，Lipschitz 常数均为 $1$。此外，\(T\)还是$\mathbb{R}^n$上的\(C^\infty\) 微分同胚.
\end{lemma}
\begin{proof}

对任意 $y_1,y_2\in\mathbb{R}^n$，由 $Q^T$ 为正交矩阵且正交矩阵保持范数可知
\begin{align*}
|T(y_1)-T(y_2)|
=|Q^T(y_1-y_2)|
=|y_1-y_2|.
\end{align*}
故 $T$ 为等距映射，特别地 $T$ 是单射。

对任意 $x\in\mathbb{R}^n$，令
$y\triangleq Q(x-x_0),$
则
\begin{align*}
T(y)=x_0+Q^T Q(x-x_0)=x_0+(x-x_0)=x.
\end{align*}
故 $T$ 是满射。因此 $T$ 是双射，且其逆映射为
\begin{align*}
T^{-1}(x)=Q(x-x_0).
\end{align*}
对任意 $x_1,x_2\in\mathbb{R}^n$，同样有
\begin{align*}
|T^{-1}(x_1)-T^{-1}(x_2)|
=|Q(x_1-x_2)|
=|x_1-x_2|.
\end{align*}
所以 $T$ 与 $T^{-1}$ 都是 Lipschitz 映射，Lipschitz 常数均为 $1$。

由上述证明可知
\begin{align*}
T(y)=x_0+Q^T y,\qquad T^{-1}(x)=Q(x-x_0),\qquad y,x\in\mathbb R^n.
\end{align*}
对仿射映射求导得
\begin{align*}
DT(y)=Q^T,\qquad DT^{-1}(x)=Q,\qquad y,x\in\mathbb R^n.
\end{align*}
右端均为常数矩阵, 故 \(DT\) 与 \(DT^{-1}\) 连续, 从而 \(T,T^{-1}\in C^\infty(\mathbb R^n)\). 又由上述证明知 \(T\) 为双射, 因此 \(T\) 是$\mathbb{R}^n$上的\(C^\infty\) 微分同胚.
\end{proof}

\begin{theorem}[直到边界光滑的函数整体逼近定理]\label{theorem:ch05-直到边界光滑的函数整体逼近定理}
设 $U\subseteq\mathbb{R}^{n}$ 是有界开集, $\partial U$ 为 $C^{1}$, 且 $u\in W^{k,p}(U)$, 其中 $k\in\mathbb{N}$, $1\leqslant p<\infty$. 则存在函数列 $u_m\in C^{\infty}(\overline{U})$, 使得
\begin{align*}
\lim\limits_{m\to\infty}\|u_m-u\|_{W^{k,p}(U)}=0.
\end{align*}
若还满足 $u\in C(\overline{U})$, 则上述函数列可以选得同时满足
\begin{align*}
\lim\limits_{m\to\infty}\sup\limits_{x\in\overline{U}}|u_m(x)-u(x)|=0,
\end{align*}
即此时$u_m$在$\overline{U}$上一致收敛到$u$.
\end{theorem}
\begin{remark}
局部构造的关键是先将取值点沿指向区域内部的方向平移, 再进行磨光. 平移距离必须足以保证整个磨光球包含在 $U$ 内, 同时还要保证磨光球没有离开描述边界的坐标邻域.
\end{remark}
\begin{proof}
\noindent$\mathbf{Step}\bm{1}$.{\heiti 确定边界坐标变换及平移方向.}

固定 $x_0\in\partial U$. 根据\cref{definition:appC-Ck边界}, 存在 $r>0$, 正交矩阵 $Q$ 及函数 $\gamma\in C^{1}(\mathbb{R}^{n-1})$, 使得在坐标变换
\begin{align*}
T(x)\triangleq Q(x-x_0),\quad
\widetilde{U}\triangleq T(U)
\end{align*}
下, 有
\begin{align*}
\widetilde{U}\cap B(0,r)
=
\{y=(y',y_n)=(y_1,\cdots,y_{n-1},y_n)\in B(0,r)\mid y_n>\gamma(y')\},
\end{align*}
其中$Q$ 的最后一行为 $-\nu(x_0)^{T}$,$\nu(x_0)$ 是$\partial U$在$x_0$处的单位外法向量.
由\hyperref[lemma:正交平移变换为等距双射]{$T$是$C^\infty$微分同胚}和\cref{Topology-theorem:同胚保持集合内点闭包边界}知$T(x_0)=0\in\partial\widetilde{U}$.又因为$\widetilde{U}$ 是开集, 而 $0\in\partial\widetilde{U}$, 因此 $0\notin\widetilde{U}$. 根据
\begin{align*}
\widetilde{U}\cap B(0,r)
=
\{y\in B(0,r)\mid y_n>\gamma(y')\},
\end{align*}
如果 $\gamma(0)<0$, 则原点满足 $0>\gamma(0)$, 从而 $0\in\widetilde{U}$, 矛盾! 故
\begin{align*}
\gamma(0)\geqslant0.
\end{align*}
另一方面, 由 $0\in\partial\widetilde{U}$, 可以选取点列
$y^{(j)}\in\widetilde{U}\cap B(0,r)$, 使得 $y^{(j)}\to0$. 对每个 $j$, 都有
\begin{align*}
y_n^{(j)}>\gamma\bigl((y^{(j)})'\bigr).
\end{align*}
令 $j\to\infty$, 利用 $\gamma$ 的连续性可得
\begin{align*}
0\geqslant\gamma(0).
\end{align*}
综合两方面, 得到 $\gamma(0)=0$. 
令
\begin{align*}
a\triangleq Q^{T}e_n.
\end{align*}
由 $Q^{T}Q=QQ^{T}=I$ 可知
\begin{align*}
|a|=1,\quad
T(x+ta)=T(x)+te_n.
\end{align*}
因此, 新坐标中沿 $e_n$ 的正方向平移, 对应原坐标中沿 $a$ 的正方向平移. 由于区域在边界图像的上方, 这正是局部指向 $U$ 内部的方向.

记
\begin{align*}
V\triangleq U\cap B\left(x_0,\frac{r}{2}\right),\quad
K\triangleq\overline{U}\cap\overline{B\left(x_0,\frac{r}{2}\right)}.
\end{align*}
取
\begin{align*}
L\triangleq\max\limits_{|y'|\leqslant r}|\nabla\gamma(y')|,\quad
\lambda\triangleq L+2,\quad
\varepsilon_{*}\triangleq\frac{r}{4(\lambda+1)}.
\end{align*}
由\hyperref[Basis of Analytics-theorem:多元函数带Peano余项的Taylor公式]{多元函数微分中值定理}, 对 $|y'|,|z'|\leqslant r$, 有
\begin{align}\label{eq::JIOJAMZKZMCBNVZXVV1}
|\gamma(y')-\gamma(z')|\leqslant L|y'-z'|.
\end{align}
当 $n=1$ 时,取 $L=0$.

下面证明, 对任意 $0<\varepsilon<\varepsilon_{*}$ 及任意 $x\in K$, 都有
\begin{align}\label{eq::JIOJAMZKZMCBNVZXVV}
\overline{B(x+\lambda\varepsilon a,\varepsilon)}
\subseteq U\cap B(x_0,r).
\end{align}
事实上, 任取 $|z|\leqslant\varepsilon$, 令
\begin{align*}
y\triangleq T(x),\quad
q\triangleq T(x+\lambda\varepsilon a-z)
=y+\lambda\varepsilon e_n-Qz.
\end{align*}
由于 $x\in K$, 有 $|y|\leqslant\frac{r}{2}$ 及 $y_n\geqslant\gamma(y')$. 又因为 $Q$ 保持欧氏长度, 所以
\begin{align*}
|q'|\leqslant |q|
\leqslant |y|+\lambda\varepsilon+|Qz|\leqslant\frac{r}{2}+(\lambda+1)\varepsilon
<\frac{3r}{4}<r.
\end{align*}
再由\eqref{eq::JIOJAMZKZMCBNVZXVV1}式可得
\begin{align*}
&q_n-\gamma \bigl( q^{\prime} \bigr) =y_n+\lambda \varepsilon -\left( Qz \right) _n-\gamma \bigl( q^{\prime} \bigr) 
\\
&=y_n+\lambda \varepsilon -\left| \left( Qz \right) _n \right|-\left[ \gamma \bigl( q^{\prime} \bigr) -\gamma \left( y^{\prime} \right) \right] -\gamma \left( y^{\prime} \right) 
\\
&\geqslant y_n-\gamma \left( y^{\prime} \right) +\lambda \varepsilon -\left| \left( Qz \right) _n \right|-L\left| q^{\prime}-y^{\prime} \right|
\\
&=y_n-\gamma (y^{\prime})+\lambda \varepsilon -\left| \left( Qz \right) _n \right|-L\left| \left( Qz \right) ^{\prime} \right|
\\
&\geqslant \lambda \varepsilon -\left( 1+L \right) \left| z \right|\geqslant \left( \lambda -1-L \right) \varepsilon \geqslant \varepsilon >0.
\end{align*}
故 $q\in\widetilde{U}\cap B(0,r)$, 即
$x+\lambda\varepsilon a-z\in U\cap B(x_0,r)$, 故\eqref{eq::JIOJAMZKZMCBNVZXVV}式得证.见\cref{figure:ch05-wiq12}.
\begin{figure}[H]
\centering
\includegraphics[width=0.5\textwidth]{ch05-wiq12.png}
\caption{}
\label{figure:ch05-wiq12}
\end{figure}

\noindent$\mathbf{Step}\bm{2}$.{\heiti 完整定义平移函数$\{u_\varepsilon\}$和磨光函数$\{v_\varepsilon\}$.}

固定非负磨光核 $\eta\in C_c^{\infty}(B(0,1))$, 满足
\begin{align*}
\int_{\mathbb{R}^{n}}\eta(z)\,\mathrm{d}z=1,\quad
\eta_{\varepsilon}(z)
\triangleq\frac{1}{\varepsilon^{n}}
\eta\left(\frac{z}{\varepsilon}\right).
\end{align*}
对每个满足 $|\alpha|\leqslant k$ 的多重指标 $\alpha$, 将 $D^{\alpha}u$ 分别作零延拓, 记为
\begin{align*}
F_{\alpha}(x)\triangleq
\begin{cases}
D^{\alpha}u(x),&x\in U,\\
0,&x\in\mathbb{R}^{n}\setminus U.
\end{cases}
\end{align*}
于是 $F_{\alpha}\in L^{p}(\mathbb{R}^{n})$, 其中 $F_0$ 是 $u$ 的零延拓. 
对 $0<\varepsilon<\varepsilon_{*}$, 在整个 $\mathbb{R}^{n}$ 上定义
\begin{align*}
u_{\varepsilon}(x)
\triangleq F_0(x+\lambda\varepsilon a),\quad
v_{\varepsilon}(x)
\triangleq(\eta_{\varepsilon}*u_{\varepsilon})(x)
=\int_{\mathbb{R}^{n}}
\eta_{\varepsilon}(z)
F_0(x+\lambda\varepsilon a-z)\,\mathrm{d}z.
\end{align*}
由\hyperref[theorem:appC-光滑算子的性质]{光滑算子的性质\ref{theorem:appC-光滑算子的性质-1}}, 有
\begin{align}
v_{\varepsilon}\in C^{\infty}(\mathbb{R}^{n}).\label{dsjfiosjaiojio2jioewj}
\end{align}

令
\begin{align*}
G_{\varepsilon}
\triangleq\{x\in\mathbb{R}^{n}\mid x+\lambda\varepsilon a\in U\}.
\end{align*}
对任意
$\varphi\in C_c^{\infty}(G_{\varepsilon})$,记$\psi(y)\triangleq \phi(y-\lambda\varepsilon a)$,则$\psi \in C_c^\infty(U)$.对 $|\alpha|\leqslant k$,换元并使用$u$的弱导数定义,有
\begin{align*}
\int_{G_{\varepsilon}}
u_{\varepsilon}(x)D^{\alpha}\varphi(x)\,\mathrm{d}x
&=\int_U
u(y)D^{\alpha}\varphi(y-\lambda\varepsilon a)\,\mathrm{d}y\\
&=(-1)^{|\alpha|}
\int_U
D^{\alpha}u(y)\varphi(y-\lambda\varepsilon a)\,\mathrm{d}y\\
&=(-1)^{|\alpha|}
\int_{G_{\varepsilon}}
D^{\alpha}u(x+\lambda\varepsilon a)
\varphi(x)\,\mathrm{d}x.
\end{align*}
因此
\begin{align}\label{eq::NIOUHNANNZZZMVBZMNX}
D^{\alpha}u_{\varepsilon}(x)
=(D^{\alpha}u)(x+\lambda\varepsilon a).
\end{align}
再由$u\in W^{k,p}(U)$,换元可得
\begin{align*}
\int_{G_{\varepsilon}}
|D^{\alpha}u_{\varepsilon}(x)|^p\,\mathrm{d}x
=
\int_{G_{\varepsilon}}
|(D^{\alpha}u)(x+\lambda\varepsilon a)|^p\,\mathrm{d}x=
\int_U |D^{\alpha}u(y)|^p\,\mathrm{d}y
<\infty.
\end{align*}
故$u_{\varepsilon}\in W^{k,p}(G_{\varepsilon}).$

由\eqref{eq::JIOJAMZKZMCBNVZXVV}式, 对每个 $x\in K$, 函数$g\left( y \right) \triangleq \eta _{\varepsilon}(x+\lambda \varepsilon a-y)$的支集满足
\begin{align*}
\mathrm{spt}\,g&\subseteq \left\{ y\in \mathbb{R} ^n:x+\lambda \varepsilon a-y\in \overline{B\left( 0,\varepsilon \right) } \right\} =\left\{ y\in \mathbb{R} ^n:y\in \overline{B\left( x+\lambda \varepsilon a,\varepsilon \right) } \right\} 
\\
&=\overline{B\left( x+\lambda \varepsilon a,\varepsilon \right) }\subseteq U\cap B\left( x_0,r \right) \subseteq U,
\end{align*}
即$g\in C_c^\infty(U)$.因而利用$u$的弱导数定义, 结合\hyperref[theorem:appC-光滑算子的性质]{光滑算子的性质\ref{theorem:appC-光滑算子的性质-1}}, 对 $|\alpha|\leqslant k$ 有
\begin{align}\label{NJHIOUSNKNZKHJFUH3251}
D^{\alpha}v_{\varepsilon}(x)
=\int_U
D_x^{\alpha}\eta_{\varepsilon}(x+\lambda\varepsilon a-y)
u(y)\,\mathrm{d}y=\int_U
\eta_{\varepsilon}(x+\lambda\varepsilon a-y)
D^{\alpha}u(y)\,\mathrm{d}y=\int_{\mathbb{R}^{n}}
\eta_{\varepsilon}(z)
F_{\alpha}(x+\lambda\varepsilon a-z)\,\mathrm{d}z.
\end{align}

\noindent$\mathbf{Step}\bm{3}$.{\heiti 证明$\{v_\varepsilon\}$局部Sobolev收敛到$u$.}

对 $t>0$, 定义
\begin{align*}
\omega_{\alpha}(t)
\triangleq
\sup\limits_{|h|\leqslant t}
\|F_{\alpha}(\cdot+h)-F_{\alpha}\|_{L^{p}(\mathbb{R}^{n})},\quad
\Omega(t)
\triangleq
\sum\limits_{|\alpha|\leqslant k}\omega_{\alpha}(t).
\end{align*}
由于 $1\leqslant p<\infty$, 根据\hyperref[Real Analysis-theorem:Lp范数下平均连续性]{$L^{p}$ 范数的平均连续性}, 有
\begin{align}\label{JGIONOMKZLMNMNZV}
\lim\limits_{t\to 0}\omega_{\alpha}(t)=0,\quad
\lim\limits_{t\to 0}\Omega(t)=0.
\end{align}

对任意 $|\alpha|\leqslant k$, 先用三角不等式分解为
\begin{align}
\|D^{\alpha}v_{\varepsilon}-D^{\alpha}u\|_{L^{p}(V)}
\leqslant
\|D^{\alpha}v_{\varepsilon}
-D^{\alpha}u_{\varepsilon}\|_{L^{p}(V)}
+
\|D^{\alpha}u_{\varepsilon}
-D^{\alpha}u\|_{L^{p}(V)}.
\label{align:ch05-边界逼近收敛1}
\end{align}
对\eqref{align:ch05-边界逼近收敛1}式第二项, 由\eqref{eq::NIOUHNANNZZZMVBZMNX}式可得
\begin{align}\label{eqJIOJFIOIWIOURIOWORWPPPPP223}
\|D^{\alpha}u_{\varepsilon}-D^{\alpha}u\|_{L^{p}(V)}
=
\|F_{\alpha}(\cdot+\lambda\varepsilon a)
-F_{\alpha}\|_{L^{p}(V)}\leqslant\omega_{\alpha}(\lambda\varepsilon).
\end{align}
对\eqref{align:ch05-边界逼近收敛1}式第一项, 由\eqref{NJHIOUSNKNZKHJFUH3251}式和磨光核的积分为 $1$, 得到
\begin{align*}
D^{\alpha}v_{\varepsilon}(x)
-D^{\alpha}u_{\varepsilon}(x)
=
\int_{\mathbb{R}^{n}}\eta_{\varepsilon}(z)
\bigl[
F_{\alpha}(x+\lambda\varepsilon a-z)
-F_{\alpha}(x+\lambda\varepsilon a)
\bigr]\,\mathrm{d}z
\end{align*}
对几乎处处的 $x\in V$ 成立. 应用\hyperref[Basis of Analytics-theorem:Minkowski(闵可夫斯基)不等式]{Minkowski不等式},有
\begin{align*}
\left\| D^{\alpha}v_{\varepsilon}-D^{\alpha}u_{\varepsilon} \right\| _{L^p(V)}&=\left[ \int_V{\left| \int_{\mathbb{R} ^n}{\eta _{\varepsilon}\left( z \right) \left[ F_{\alpha}\left( x+\lambda \varepsilon a-z \right) -F_{\alpha}\left( x+\lambda \varepsilon a \right) \right] \,\mathrm{d}z} \right|^p\mathrm{d}x} \right] ^{\frac{1}{p}}
\\
&\leqslant \int_{\mathbb{R} ^n}{\left( \int_V{\left| \eta _{\varepsilon}\left( z \right) \right|^p\left| F_{\alpha}\left( x+\lambda \varepsilon a-z \right) -F_{\alpha}\left( x+\lambda \varepsilon a \right) \right|^p\,\mathrm{d}x} \right) ^{\frac{1}{p}}\mathrm{d}z}
\\
&=\int_{B\left( 0,\varepsilon \right)}{\eta _{\varepsilon}\left( z \right) \left\| F_{\alpha}\left( \cdot +\lambda \varepsilon a-z \right) -F_{\alpha}\left( \cdot +\lambda \varepsilon a \right) \right\| _{L^p\left( V \right)}\,\mathrm{d}z}
\\
&\leqslant \int_{B\left( 0,\varepsilon \right)}{\eta _{\varepsilon}(z)\left\| F_{\alpha}\left( \cdot -z \right) -F_{\alpha} \right\| _{L^p(\mathbb{R} ^n)}\,\mathrm{d}z}
\\
&\leqslant \omega _{\alpha}\left( \varepsilon \right) \int_{B\left( 0,\varepsilon \right)}{\eta _{\varepsilon}(z)\,\mathrm{d}z}=\omega _{\alpha}\left( \varepsilon \right) .
\end{align*}
将上述估计和\eqref{eqJIOJFIOIWIOURIOWORWPPPPP223}式代入\eqref{align:ch05-边界逼近收敛1}式, 得到
\begin{align}\label{NJHIOUSNKNZKHJFUH325}
\|v_{\varepsilon}-u\|_{W^{k,p}(V)}
\leqslant
\sum\limits_{|\alpha|\leqslant k}
\|D^{\alpha}v_{\varepsilon}-D^{\alpha}u\|_{L^{p}(V)}\leqslant
\Omega(\varepsilon)+\Omega(\lambda\varepsilon)
\leqslant
2\Omega\bigl((\lambda+1)\varepsilon\bigr).
\end{align}
再结合\eqref{JGIONOMKZLMNMNZV}式可得
\begin{align}
\lim\limits_{\varepsilon\to 0}
\|v_{\varepsilon}-u\|_{W^{k,p}(V)}=0.
\label{align:ch05-边界逼近收敛}
\end{align}

\noindent$\mathbf{Step}\bm{4}$.{\heiti 固定有限覆盖及光滑单位分解}

由 $\partial U$ 的紧性, 可以选取有限个边界点 $x_i^0$ 及相应的半径 $r_i$, 使得
\begin{align*}
\partial U\subseteq
\bigcup\limits_{i=1}^{N}O_i,\quad
O_i\triangleq B\left(x_i^0,\frac{r_i}{2}\right).
\end{align*}
记
\begin{align*}
V_i\triangleq U\cap O_i,\quad
K_i\triangleq\overline{U}\cap\overline{O_i},
\quad i=1,\cdots,N.
\end{align*}
将各坐标片中按$\mathbf{Step1}$的方法确定的方向和常数分别记为
\begin{align*}
a_i=Q_i^{T}e_n,\quad
\lambda_i=L_i+2,\quad
\varepsilon_{i,*}=\frac{r_i}{4(\lambda_i+1)}.
\end{align*}

集合
\begin{align*}
K_0
\triangleq
\overline{U}\setminus\bigcup\limits_{i=1}^{N}O_i
\end{align*}
是 $U$ 的紧子集. 因而由\cref{proposition:紧集的相对紧开邻域},可选取开集 $V_0\triangleq O_0$, 使得
\begin{align*}
K_0\subseteq O_0,\quad
\overline{O_0}\subseteq U.
\end{align*}
于是
\begin{align*}
\overline{U}\subseteq
\bigcup\limits_{i=0}^{N}O_i.
\end{align*}

根据\hyperref[lemma:ch05-光滑单位分解-p7R3]{光滑单位分解引理}, 可选取函数
$\zeta_0,\cdots,\zeta_N$, 满足
\begin{gather*}
\zeta_i\in C_c^{\infty}(O_i),\quad
0\leqslant\zeta_i\leqslant1,\quad
i=0,\cdots,N;\\
\sum\limits_{i=0}^{N}\zeta_i(x)=1,\quad x\in \overline{U}.
\end{gather*}

\noindent$\mathbf{Step}\bm{5}$.{\heiti 构造局部磨光函数$\{v_i\}_{i=0}^{N}$.}

任取 $\delta>0$. 由\eqref{JGIONOMKZLMNMNZV}式, 可以选取 $t_{\delta}>0$, 使得
\begin{align*}
2\Omega(t_{\delta})<\delta.
\end{align*}
对每个边界坐标片,取
\begin{align*}
\varepsilon_i
\triangleq
\min\left\{
\frac{r_i}{8(\lambda_i+1)},
\frac{t_{\delta}}{2(\lambda_i+1)}
\right\},
\quad i=1,\cdots,N.
\end{align*}
于是
\begin{align*}
0<\varepsilon_i<\varepsilon_{i,*},\quad
(\lambda_i+1)\varepsilon_i
\leqslant\frac{t_{\delta}}{2}.
\end{align*}
定义
\begin{align*}
v_i(x)
\triangleq
(\eta_{\varepsilon_i}*F_0)
(x+\lambda_i\varepsilon_i a_i),
\quad x\in\mathbb{R}^{n}.
\end{align*}
由\eqref{dsjfiosjaiojio2jioewj}\eqref{NJHIOUSNKNZKHJFUH325}式知,
$v_i\in C^{\infty}(\mathbb{R}^{n})$, 且
\begin{align}
\|v_i-u\|_{W^{k,p}(V_i)}
\leqslant
2\Omega\bigl((\lambda_i+1)\varepsilon_i\bigr)
<\delta,
\quad i=1,\cdots,N.
\label{align:ch05-边界局部误差}
\end{align}

对于内部开集 $V_0$, 由于 $\overline{V_0}$ 是 $U$ 的紧子集, 并且由\subcref{Real Analysis-proposition:点集间的距离的性质}{Real Analysis-proposition:点集间的距离的性质-3}知
\begin{align*}
d_0\triangleq
\operatorname{dist}(\overline{V_0},\mathbb{R}^{n}\setminus U)>0.
\end{align*}
由\hyperref[theorem:ch05-光滑函数的局部逼近定理]{光滑函数的局部逼近定理}, 对任意 $\delta>0$, 存在$\varepsilon_0\in (0,d_0)$及函数 $v_0\triangleq v_{\varepsilon_0}\in C^\infty(\overline{V_0})$, 使得
\begin{align}
\|v_{\varepsilon_0}-u\|_{W^{k,p}(V_0)}<\delta.
\label{align:ch05-内部逼近误差}
\end{align}

\noindent$\mathbf{Step}\mathbf{6}.${\heiti 拼接函数列$\{v_i\}_{i=0}^{N}$和对应单位分解$\{\zeta_i\}_{i=0}^{N}$得到$v$并证明其光滑性与收敛性.}

由于 $V_0=O_0$, 有 $v_0\in C^\infty(O_0)$; 对 $i=1,\cdots,N$, 已有 $v_i\in C^\infty(\mathbb{R}^{n})$, 因而也有 $v_i\in C^\infty(O_i)$. 所以对每个 $i=0,\cdots,N$, 乘积 $\zeta_i v_i$ 都在 $O_i$ 内光滑.

由 $\zeta_i\in C_c^\infty(O_i)$, $\operatorname{spt}\zeta_i$ 是包含在 $O_i$ 内的紧集. 在整个 $\mathbb{R}^{n}$ 上定义
\begin{align*}
g_i(x)\triangleq
\begin{cases}
\zeta_i(x)v_i(x),&x\in O_i,\\
0,&x\in\mathbb{R}^{n}\setminus O_i.
\end{cases}
\end{align*}
下面证明 $g_i\in C_c^\infty(\mathbb{R}^{n})$.

若 $x\in O_i$, 则由于 $O_i$ 是开集, 可以选取 $r_x>0$, 使得 $B(x,r_x)\subseteq O_i$. 在这个邻域内,
\begin{align*}
g_i=\zeta_i v_i,
\end{align*}
所以 $g_i$ 在 $x$ 的某个邻域内光滑.

若 $x\notin O_i$, 则由 $\operatorname{spt}\zeta_i\subseteq O_i$ 可知 $x\notin \operatorname{spt}\zeta_i$. 因为 $\operatorname{spt}\zeta_i$ 是闭集, 可以选取 $r_x>0$, 使得
\begin{align*}
B(x,r_x)\cap \operatorname{spt}\zeta_i=\varnothing.
\end{align*}
任取 $y\in B(x,r_x)$,则由上式知 $y\notin\operatorname{spt}\zeta_i$ 可知 $\zeta_i(y)=0$, 因而 $g_i(y)=0$.因此
\begin{align*}
g_i\equiv0\quad\text{在 }B(x,r_x)\text{ 上}.
\end{align*}
所以 $g_i$ 在 $x$ 的某个邻域内也是光滑的. 

综上, $g_i$ 在 $\mathbb{R}^{n}$ 的每一点附近都光滑, 故$g_i\in C^\infty(\mathbb{R}^{n}).$
又因为 $g_i$ 在 $\mathbb{R}^{n}\setminus \operatorname{spt}\zeta_i$ 上恒为零, 且 $\operatorname{spt}\zeta_i$ 为闭集, 所以$\operatorname{spt}g_i\subseteq \operatorname{spt}\zeta_i.$
由 $\operatorname{spt}\zeta_i$ 的紧性可得 $g_i\in C_c^\infty(\mathbb{R}^{n})$.

现在定义
\begin{align*}
v^{(\delta)}\triangleq\sum\limits_{i=0}^{N}g_i.
\end{align*}
这是有限个全空间光滑函数之和, 因而
$v^{(\delta)}\in C^\infty(\mathbb{R}^{n}).$
并且
\begin{align*}
\operatorname{spt}v^{(\delta)}
\subseteq
\bigcup\limits_{i=0}^{N}\operatorname{spt}g_i
\subseteq
\bigcup\limits_{i=0}^{N}\operatorname{spt}\zeta_i.
\end{align*}
右端是有限个紧集之并, 仍为紧集, 故
\begin{align*}
v^{(\delta)}\in C_c^\infty(\mathbb{R}^{n}).
\end{align*}
因此
\begin{align*}
v^{(\delta)}|_{\overline{U}}\in C^\infty(\overline{U}).
\end{align*}

最后,因为在 $U$ 上 $\sum\limits_{i=0}^{N}\zeta_i=1$, 所以
\begin{align*}
v^{(\delta)}(x)-u(x)=\sum\limits_{i=0}^{N}\zeta_i(x)(v_i(x)-u(x)),\quad \forall x\in U.
\end{align*}
根据\hyperref[theorem:ch05-弱导数性质]{弱导数性质\ref{theorem:ch05-弱导数性质-4}}, 又由$\zeta_i\in C_c^{\infty}(O_i)$知
$\operatorname{spt}D^{\beta}\zeta_i\subseteq O_i$, 再利用\eqref{align:ch05-边界局部误差}\eqref{align:ch05-内部逼近误差}式有
\begin{align*}
\|v^{(\delta)}-u\|_{W^{k,p}(U)}
&\leqslant
\sum\limits_{|\alpha|\leqslant k}
\|D^{\alpha}(v^{(\delta)}-u)\|_{L^{p}(U)}\leqslant
\sum\limits_{i=0}^{N}
\sum\limits_{|\alpha|\leqslant k}
\sum\limits_{\beta\leqslant\alpha}
\binom{\alpha}{\beta}
\|D^{\beta}\zeta_i\|_{L^{\infty}(\mathbb{R}^{n})}
\|D^{\alpha-\beta}(v_i-u)\|_{L^{p}(V_i)}\\
&\leqslant
\sum\limits_{i=0}^{N}
C_i\|v_i-u\|_{W^{k,p}(V_i)}\leqslant
\left(\sum\limits_{i=0}^{N}C_i\right)\delta
=C\delta ,
\end{align*}
其中\begin{align*}
C_i
\triangleq
\sum\limits_{|\alpha|\leqslant k}
\sum\limits_{\beta\leqslant\alpha}
\binom{\alpha}{\beta}
\|D^{\beta}\zeta_i\|_{L^{\infty}(\mathbb{R}^{n})},\quad
C\triangleq \sum\limits_{i=0}^{N}C_i.
\end{align*}

综上,对每个正整数 $m$, 取$\delta_m\triangleq\frac{1}{mC},$
根据$\mathbf{Step}\mathbf{1}-\mathbf{Step}\mathbf{6}$相应构造出在$\mathbb{R}^n$上的光滑函数$v^{(\delta_m)}$,并记$u_m\triangleq v^{(\delta_m)}|_{\overline{U}}$,则
\begin{align*}
u_m\in C^{\infty}(\overline{U}),\quad
\|u_m-u\|_{W^{k,p}(U)}<\frac{1}{m}.
\end{align*}

\noindent$\mathbf{Step}\mathbf{7}$.{\heiti 连续情形下同时保证一致收敛.}

进一步假设 $u\in C(\overline{U})$. 
任取 $\delta>0$. 由于 $\overline{U}$ 紧, $u$ 在 $\overline{U}$ 上一致连续. 因此存在 $\rho_\delta>0$, 使得对任意 $x,y\in\overline{U}$, 都有
\begin{align}\label{eqL::JIONSCIOUSAJiow}
|x-y|<\rho_\delta
\Longrightarrow
|u(x)-u(y)|<\frac{\delta}{2}.
\end{align}
另一方面, 由 $\mathbf{Step}\mathbf{3}$ 中的\eqref{JGIONOMKZLMNMNZV}式, 可以选取 $t_\delta>0$, 同时满足
\begin{align*}
t_\delta<\rho_\delta,\quad
2\Omega(t_\delta)<\delta.
\end{align*}

(i)先处理边界坐标片. 对 $i=1,\cdots,N$, 仍按 $\mathbf{Step}\mathbf{5}$ 的方法取
\begin{align*}
\varepsilon_i
\triangleq
\min\left\{
\frac{r_i}{8(\lambda_i+1)},
\frac{t_\delta}{2(\lambda_i+1)}
\right\},
\end{align*}
并令
\begin{align*}
v_i(x)
\triangleq
(\eta_{\varepsilon_i}*F_0)
(x+\lambda_i\varepsilon_i a_i),
\quad x\in\mathbb{R}^{n}.
\end{align*}
于是
\begin{align*}
0<\varepsilon_i<\varepsilon_{i,*},\quad
(\lambda_i+1)\varepsilon_i
\leqslant\frac{t_\delta}{2}
<\rho_\delta.
\end{align*}
从而由 $\mathbf{Step}\mathbf{3}$ 中的\eqref{NJHIOUSNKNZKHJFUH325}式可得
\begin{align*}
\|v_i-u\|_{W^{k,p}(V_i)}
\leqslant
2\Omega\bigl((\lambda_i+1)\varepsilon_i\bigr)
\leqslant2\Omega(t_\delta)
<\delta.
\end{align*}

(ii)下面估计边界坐标片上的一致误差. 任取
\begin{align*}
x\in K_i=\overline{U}\cap\overline{O_i}.
\end{align*}
由 $\mathbf{Step}\mathbf{1}$ 中的\eqref{eq::JIOJAMZKZMCBNVZXVV}式, 对任意 $z\in B(0,\varepsilon_i)$, 都有
\begin{align*}
x+\lambda_i\varepsilon_i a_i-z\in U,
\end{align*}
并且
\begin{align*}
|x+\lambda_i\varepsilon_i a_i-z-x|
\leqslant\lambda_i\varepsilon_i+|z|
\leqslant(\lambda_i+1)\varepsilon_i
<\rho_\delta.
\end{align*}
因此, 由\eqref{eqL::JIONSCIOUSAJiow}式可知
\begin{align*}
|u(x+\lambda_i\varepsilon_i a_i-z)-u(x)|
<\frac{\delta}{2}.
\end{align*}
又因为 $\eta_{\varepsilon_i}\geqslant0$, 支撑包含在 $B(0,\varepsilon_i)$ 内, 且积分为 $1$, 所以
\begin{align*}
|v_i(x)-u(x)|
&=
\left|
\int_{B(0,\varepsilon_i)}
\eta_{\varepsilon_i}(z)
\bigl[u(x+\lambda_i\varepsilon_i a_i-z)-u(x)\bigr]
\,\mathrm{d}z
\right|\\
&\leqslant
\int_{B(0,\varepsilon_i)}
\eta_{\varepsilon_i}(z)
|u(x+\lambda_i\varepsilon_i a_i-z)-u(x)|
\,\mathrm{d}z\\
&\leqslant
\frac{\delta}{2}
\int_{B(0,\varepsilon_i)}
\eta_{\varepsilon_i}(z)\,\mathrm{d}z
=\frac{\delta}{2}.
\end{align*}
由于这个估计对所有 $x\in K_i$ 都成立,故
\begin{align}\label{1242335346345363451}
\sup\limits_{x\in K_i}|v_i(x)-u(x)|
\leqslant\frac{\delta}{2},
\quad i=1,\cdots,N.
\end{align}

(iii)再处理内部开集 $V_0$. 
由于 $\overline{V_0}$ 是 $U$ 的紧子集,并且由\subcref{Real Analysis-proposition:点集间的距离的性质}{Real Analysis-proposition:点集间的距离的性质-3}知
\begin{align*}
d_0\triangleq
\operatorname{dist}(\overline{V_0},\mathbb{R}^{n}\setminus U)>0.
\end{align*}
利用\hyperref[theorem:ch05-光滑函数的局部逼近定理]{光滑函数的局部逼近定理}知,存在$\delta_1>0,$当$0<\varepsilon<
\min\left\{d_0,\delta_1\right\}$时,有
\begin{align}\label{ZNMBZNBNVJZBNMCOIOAIAKJAAA}
w_\varepsilon\triangleq \eta_\varepsilon * u\in C^\infty(\overline{V_0}),\quad \|w_\varepsilon-u\|_{W^{k,p}(V_0)}<\delta .
\end{align}
同时, 当$0<\varepsilon<
\min\left\{\frac{d_0}{2},\frac{\rho_\delta}{2}\right\}$时, 对任意 $x\in\overline{V_0}$ 及 $z\in B(0,\varepsilon)$, 都有
\begin{align*}
x-z\in U,\quad
|x-z-x|=|z|<\rho_\delta.
\end{align*}
因此,由\eqref{eqL::JIONSCIOUSAJiow}式可得
\begin{align*}
|u(x-z)-u(x)|<\frac{\delta}{2}.
\end{align*}
又因为 $\eta_{\varepsilon}\geqslant0$, 支撑包含在 $B(0,\varepsilon_i)$ 内, 且积分为 $1$, 所以对任意 $x\in\overline{V_0}$有
\begin{align*}
|w_\varepsilon(x)-u(x)|
\leqslant
\int_{B(0,\varepsilon)}
\eta_\varepsilon(z)|u(x-z)-u(x)|\,\mathrm{d}z
\leqslant
\frac{\delta}{2}
\int_{B(0,\varepsilon)}
\eta_\varepsilon(z)\,\mathrm{d}z
=\frac{\delta}{2}.
\end{align*}
从而
\begin{align}\label{ZNMBZNBNVJZBNMCOIOAIAKJAAA---2}
\sup\limits_{x\in\overline{V_0}}
|w_\varepsilon(x)-u(x)|
\leqslant\frac{\delta}{2}.
\end{align}
结合\eqref{ZNMBZNBNVJZBNMCOIOAIAKJAAA}\eqref{ZNMBZNBNVJZBNMCOIOAIAKJAAA---2}式,可选取充分小的$\varepsilon _0\in \left( 0,\frac{\min \left\{ d_0,\rho _{\delta},\delta _1 \right\}}{2} \right)$,记$v_0\triangleq w_{\varepsilon_0},$则
\begin{align}\label{124233534634536345}
v_0\in C^{\infty}(\overline{V_0}),\quad \left\| v_0-u \right\| _{W^{k,p}(V_0)}<\delta ,\quad \mathop {\mathrm{sup}} \limits_{x\in \overline{V_0}}|v_0(x)-u(x)|\leqslant \frac{\delta}{2}.
\end{align}

(iv)对上述选定的 $v_0,\cdots,v_N$, 按 $\mathbf{Step}\mathbf{6}$ 定义
\begin{align}\label{HIOUJKLHIOUIOIOIOWIOIOIOIOIO11214}
g_i(x)\triangleq
\begin{cases}
\zeta_i(x)v_i(x),&x\in O_i,\\
0,&x\notin O_i,
\end{cases}
\quad i=0,\cdots,N,
\end{align}
并令
\begin{align*}
v_1^{(\delta)}
\triangleq
\sum\limits_{i=0}^{N}g_i.
\end{align*}
由 $\mathbf{Step}\mathbf{6}$ 已经证明的光滑性结论及范数估计, 得到
\begin{align*}
v_1^{(\delta)}\in C_c^\infty(\mathbb{R}^{n}),\quad
\|v_1^{(\delta)}-u\|_{W^{k,p}(U)}<C\delta.
\end{align*}

(v)现在证明整体的一致误差估计. 任取 $x\in\overline{U}$.
对每个 $i=0,\cdots,N$, 若 $x\in O_i$, 则由\eqref{1242335346345363451}\eqref{124233534634536345}式,可得
\begin{align*}
|g_i(x)-\zeta_i(x)u(x)|
=
\zeta_i(x)|v_i(x)-u(x)|
\leqslant\frac{\delta}{2}\zeta_i(x).
\end{align*}
若 $x\notin O_i$, 则由 $g_i$ 的定义式\eqref{HIOUJKLHIOUIOIOIOWIOIOIOIOIO11214}及 $\operatorname{spt}\zeta_i\subseteq O_i$, 有
$g_i(x)=\zeta_i(x)=0$, 因而上述不等式对任意$x\in\overline{U}$都成立.

再由$\mathbf{Step}\mathbf{4}$ 中 $\sum\limits_{i=0}^{N}\zeta_i(x)=1$,对任意$x\in\overline{U}$,都有
\begin{align*}
|v_1^{(\delta)}(x)-u(x)|
=
\left|
\sum\limits_{i=0}^{N}
\bigl(g_i(x)-\zeta_i(x)u(x)\bigr)
\right|\leqslant
\sum\limits_{i=0}^{N}
|g_i(x)-\zeta_i(x)u(x)|
\leqslant
\frac{\delta}{2}
\sum\limits_{i=0}^{N}\zeta_i(x)
=\frac{\delta}{2}.
\end{align*}
故
\begin{align*}
\sup\limits_{x\in\overline{U}}
|v_1^{(\delta)}(x)-u(x)|
\leqslant\frac{\delta}{2}.
\end{align*}

最后, 对每个正整数 $m$, 取
\begin{align*}
\delta_m\triangleq\frac{1}{mC},\quad
u_m\triangleq v_1^{(\delta_m)}|_{\overline{U}}.
\end{align*}
由于 $v_1^{(\delta_m)}$ 在整个 $\mathbb{R}^{n}$ 上光滑, 有
$u_m\in C^\infty(\overline{U})$. 同时还有
\begin{align*}
\|u_m-u\|_{W^{k,p}(U)}
&<C\delta_m=\frac{1}{m},\\
\sup\limits_{x\in\overline{U}}|u_m(x)-u(x)|
&\leqslant\frac{\delta_m}{2}
=\frac{1}{2mC}.
\end{align*}
因此, 所选取的同一函数列既在 $W^{k,p}(U)$ 中收敛于 $u$, 又在 $\overline{U}$ 上一致收敛于 $u$.
\end{proof}










\end{document}
